% !TEX root = ../main.tex

\chapter{多源全身运动数据采集与重定向}

\section{系统架构概述}

为了支撑人形机器人的全身运动控制与任务级行为学习，需要采集机器人本身的动作数据以及与环境交互的数据，提供类人运动的先验知识，以便于进行强化学习、模仿学习等。
而机器人本体运动数据又很难获得，人类动作数据虽然容易获取，但人体与机器人的骨骼结构存在差异，需要将人体动作映射为机器人可执行的关节轨迹或末端目标，此为重定向问题。

所以为了解决“数据从哪里来”和“如何稳定映射到机器人”这两个基础问题，本章构建了多源全身运动数据采集与重定向系统。
数据采集模块负责从不同设备获取操作者的姿态、关节角和轨迹信息，重定向模块负责将人体动作映射为机器人可执行的关节轨迹或末端目标。

多源数据采集的主要原因在于不同实验场景对数据精度、部署效率和
任务类型的要求差异显著。惯性动捕适合快速部署和大范围日常采集，光学动捕
能够提供更高精度的科研数据，而机器人-物体交互数据能够描述机器人与环境交互过程中的复合动作模式。


在整体流程上，操作者的人体动作首先由采集设备转换为人体姿态表示；
随后系统在统一骨架模型与坐标系下进行预处理，并根据具体机器人平台建立从人体空间到机器人空间的映射；
最后，通过基于 GMR 的重定向策略生成满足关节限幅和运动约束的机器人可执行数据。
如图~\ref{fig:gmr-framework} 所示，Noitom 惯性动捕与 Nokov 光学动捕两类人体运动数据统一汇入 GMR 模块，并进一步映射到 Unitree G1 与全尺寸机器人平台，从而形成多源数据到多机器人平台的统一重定向链路。

\begin{figure}[ht]
  \centering
  \includegraphics[width=0.85\textwidth]{Capture-GMR-G1.png}
  \caption{多源动捕数据到机器人平台的 GMR 重定向框架示意图}
  \label{fig:gmr-framework}
\end{figure}

\section{多源人体运动捕捉系统}

\subsection{惯性动捕系统}

Noitom 惯性动捕系统采用穿戴式惯性测量单元获取人体各关键骨骼段的姿态变化，具有部署方便、对场地依赖低和成本相对可控等优势。对于需要快速采集示教动作、重复采集日常操控行为或在非标准实验环境中开展测试的场景，惯性动捕系统能够提供较高的实用价值。本文将其作为常规遥操作数据采集的重要组成部分，重点服务于基础上肢动作、全身协调动作以及初步行为演示数据的积累。

尽管惯性动捕在可部署性方面表现突出，但其在长时间采集时可能存在漂移累积、局部姿态解算误差以及磁场环境敏感等问题。因此，在数据接入阶段，本文对原始关节角序列进行了统一时间戳对齐和平滑处理，并在映射前加入异常帧检测与短时插值，以提升下游重定向的稳定性。对于需要更高精度的位置级分析的实验，惯性动捕数据更多用于提供全身动作趋势与相对姿态信息。

\subsection{光学动捕系统}

Nokov 光学动捕系统通过多相机协同跟踪反光标记点，可较高精度地恢复人体关键部位的三维位置与姿态信息。与惯性动捕相比，光学动捕更加适合对末端轨迹误差、关节运动一致性和重定向精度进行定量分析，因此在需要精确评价动作还原质量和系统性能的实验中具有重要作用。

本文利用光学动捕系统采集了面向机器人操作与协调动作的高精度示教数据，并将其作为重定向算法评估的重要基准数据来源。为了便于与机器人控制接口衔接，本文将光学动捕输出进一步转换为统一骨架节点位姿表示，并对关键末端执行器轨迹进行重采样和速度约束处理。这样既能保留高精度空间信息，又能避免直接使用离散点序列时出现的轨迹抖动问题。

需要指出的是，光学动捕系统对场地布局、相机标定质量和遮挡情况较为敏感，因此其更适合用于实验室条件下的标准化采集。本文通过将光学动捕与惯性动捕相互补充，在保证高精度验证能力的同时兼顾了系统整体的采集灵活性。

\subsection{机器人-物体交互数据采集}

除人体动作采集外，本研究还基于 SONIC 框架进行机器人-物体交互数据采集。SONIC 是一个全身运动控制框架，基于其稳定、灵活的控制接口，可以方便地进行机器人-物体交互数据采集。
利用PICO基于SONIC框架进行宇树G1机器人全身遥操作、移动操作一体化数据采集，主要采集了三路数据：

1. PICO遥操作端：记录上游人体动作数据；

2. 宇树G1机器人本体：记录机器人本体运动数据；

3. 相机图像链路：记录机器人与物体交互过程中的图像信息；

为了统一说明机器人--物体交互数据采集中的多模态日志结构，表~\ref{tab:sonic-data-fields} 对三路数据中具有代表性的文件、关键字段及其用途进行了整理。该表既可作为后续数据对齐与预处理的依据，也可作为训练与评测阶段选择输入模态时的参考。

\begingroup
\footnotesize
\setlength{\LTleft}{0pt}
\setlength{\LTright}{0pt}
\renewcommand{\arraystretch}{1.1}
\begin{longtable}{@{}p{0.24\textwidth} p{0.34\textwidth} p{0.34\textwidth}@{}}
\caption{机器人--物体交互数据采集中的代表性文件与关键字段}
\label{tab:sonic-data-fields} \\
\toprule
代表文件 & 关键字段/内容 & 说明 \\
\midrule
\endfirsthead

\caption[]{机器人--物体交互数据采集中的代表性文件与关键字段（续）} \\
\toprule
代表文件 & 关键字段/内容 & 说明 \\
\midrule
\endhead

\midrule
\multicolumn{3}{r}{续下页} \\
\endfoot

\bottomrule
\endlastfoot

\multicolumn{3}{@{}l@{}}{\textbf{遥操作端(teleop)}} \\
\midrule
\path{pose_*.npz} & \path{frame_vr_position} / \path{vr_position} & 左手、右手、头/颈等关键点位置 \\
\path{pose_*.npz} & \path{frame_vr_orientation} / \path{vr_orientation} & 关键点四元数 \\
\path{pose_*.npz} & \path{body_quat_w} & 身体朝向 \\
\path{pose_*.npz} & \path{joint_pos} & teleop 侧生成的 joint 向量 \\
\path{pose_*.npz} & \path{frame_left_hand_joints} / \path{frame_right_hand_joints} & 手部角度 \\

\midrule
\multicolumn{3}{@{}l@{}}{\textbf{G1机器人本体(deploy\_logs)}} \\
\midrule
\path{q.csv} & 29D body joints & 机器人关节状态主来源 \\
\path{left_hand_q.csv} / \path{right_hand_q.csv} & 手部关节 & hand state \\
\path{base_quat.csv} & base IMU quaternion & 姿态、FK 输入 \\
\path{base_ang_vel.csv} & IMU gyro & robot state IMU \\
\path{base_accel.csv} & IMU accel & robot state IMU \\
\path{motor_torque.csv} & 29D torque & \path{tau_ff} \\
\path{torso_quat.csv} & torso quaternion & \path{torso_rpy}、\path{torso_dyaw} \\
\path{torso_ang_vel.csv} & torso angular velocity & \path{torso_vyaw}，积分得到 \path{target_yaw} \\
\path{command_velocity.csv} & \path{vx}/\path{vy}/\path{vyaw} command & planner/locomotion 命令\\
\path{command_height.csv} & torso height command & planner 高度命令\\
\path{command_heading.csv} & \path{target_yaw} / \path{delta_heading} & planner 朝向命令 \\

\midrule
\multicolumn{3}{@{}l@{}}{\textbf{相机图像链路(images)}} \\
\midrule
\path{images/index.jsonl} & \path{frame_id}、\path{timestamp_monotonic}、\path{streams} & 图像索引，不存像素本体 \\
\path{images/rgb/*.jpg} & RGB 图像 & 视觉观测 \\
\path{images/depth/*.npy} & depth & 深度观测 \\
\end{longtable}
\endgroup

从上述数据组织可以看出，遥操作端主要提供人体动作先验，部署日志主要提供机器人本体状态与控制命令，而图像链路则补充环境与目标物体的视觉观测。三者共同构成了后续全身控制学习、离线回放分析和机器人--物体交互建模的数据基础。


\section{基于 GMR 的重定向}

\subsection{GMR 算法介绍}

本文在 General Motion Retargeting（GMR）框架基础上实现人体动作到人形机器人关节空间的重定向。其核心思想并非直接复制人体关节角，而是先将不同来源的人体动作统一表示为标准化骨架帧，再结合机器人运动学模型和逆运动学任务求解机器人当前时刻的可执行位姿。通过这种方式，人体动作中的空间位姿信息能够在保留主要运动趋势的前提下映射到不同结构的人形机器人平台。

在具体实现中，GMR 的输入为单帧人体骨架数据，其数据结构由人体关节名称索引，内容包含关节位置与四元数姿态。系统首先依据人体根节点对动作数据进行尺度归一化，并根据预设的人体骨架缩放表对躯干、上下肢等关键部位进行长度修正，以减小不同操作者身高差异对重定向结果的影响。随后，系统加载目标机器人的 MuJoCo 模型以及相应的 IK 配置文件，从而获得机器人根节点定义、人体与机器人关键刚体之间的对应关系、位置偏置、姿态偏置以及各任务项的权重设置。

在求解阶段，GMR 将人体关键骨骼的目标位姿转换为机器人刚体的跟踪任务，并构建两阶段逆运动学问题。第一阶段主要保证根节点、下肢和躯干等主体结构的稳定对齐，第二阶段进一步细化上肢和末端姿态，从而在整体稳定性和局部精度之间取得平衡。每一帧重定向结果最终以机器人广义坐标 $qpos$ 的形式输出，包括根位置、根姿态以及各关节自由度角，为后续动作回放、在线控制与策略学习提供统一接口。

\subsection{基于 GMR 的多源数据重定向实现}

在本文的系统中，GMR 不仅作为一个通用重定向求解器存在，更承担了多源人体运动数据统一接入的核心作用。针对不同采集设备的数据格式差异，本文分别实现了 Noitom 动捕数据的离线重定向链路和 Nokov 光学动捕数据的实时重定向链路，使得不同来源的人体动作能够在统一框架下转换为机器人可执行轨迹。

对于 Noitom 动捕数据，本文主要采用离线 BVH 文件重定向方式。首先读取 Noitom 导出的 BVH 文件，并通过骨架解析和前向运动学计算恢复各骨骼节点在全局坐标系下的位置与姿态；随后对不同来源 BVH 文件的坐标系、尺度单位和脚部关键点定义进行统一整理，将其转换为 GMR 所要求的人体骨架帧表示；最后将处理后的数据送入重定向模块，得到目标机器人对应的根位姿和关节轨迹。该离线链路适用于动作数据集构建、示教动作回放以及重定向效果分析。

对于 Nokov 光学动捕数据，本文进一步实现了实时重定向流程。系统通过实时订阅 Nokov 输出的数据帧，读取当前时刻各人体刚体片段的位置与姿态信息，并将其在线转换为与 BVH 路径一致的标准化人体骨架格式。之后，系统在每个时刻调用 GMR 模块求解机器人当前关节配置，并将输出结果送入可视化或下游控制接口。相较于离线 BVH 路径，该方案省去了动作导出与后处理的中间环节，能够满足实时演示、在线调试和遥操作场景下的重定向需求。

综上，本文在 GMR 框架下打通了 ``离线 BVH'' 与 ``实时流数据'' 两类重定向链路，实现了 Noitom 和 Nokov 两种主要动捕来源的统一接入。这一工作保证了第 2 章数据采集层与后续机器人控制层之间的数据表示一致性，也为不同任务场景下的人体动作迁移提供了稳定基础。

\subsection{多平台适配：G1 与 Taihu 全尺寸机器人}

除多源数据接入外，本文还在同一套 GMR 主框架下完成了面向不同机器人平台的重定向适配。平台适配的关键不在于重新设计重定向求解器，而在于针对不同机器人模型建立对应的 MuJoCo 运动学模型与 IK 配置文件，使人体骨架目标能够映射到具体平台的根节点、躯干、上下肢及末端刚体上。

对于 Unitree G1 平台，本文基于其标准人形结构建立了与 GMR 相匹配的机器人模型和关节映射关系。通过在 IK 配置中定义人体骨架节点与 G1 关键刚体之间的对应关系，并设置相应的位置权重、姿态权重以及偏置参数，系统能够较稳定地完成从人体动作到 G1 关节空间的重定向。该平台构成了本文遥操作与全身动作迁移实验的主要验证对象。

在此基础上，本文进一步面向 Taihu 全尺寸机器人构建了专用的 IK 配置文件。与 G1 相比，Taihu 在机器人根节点命名、关节结构和刚体布局上均存在明显差异，因此需要重新定义人体根节点到机器人根节点的对应关系，并重新设置骨架缩放表、位置偏置、姿态偏置以及两阶段 IK 匹配表中的任务项。通过这些平台专用参数的调整，本文将同一套 GMR 重定向流程迁移到全尺寸机器人上，实现了跨平台的人体动作映射。

因此，本文的多平台适配策略可以概括为“统一重定向框架 + 平台专用配置文件”。该策略既保留了 GMR 框架在算法层面的统一性，又通过配置层的灵活调整满足了不同机器人平台的结构差异，从而降低了系统在多平台部署中的重复开发成本。
